%%%%%%%%%%%%%%%%%%%%%%%%%%%%%%%%%%%%%%%%%%%%%
\section{Introduction}
\label{Sec:Introduction}
%%%%%%%%%%%%%%%%%%%%%%%%%%%%%%%%%%%%%%%%%%%%%

Collisionless kinetic gases provide the standard description for systems consisting of a large number of particles which interact with each other only through the mean field they generate. Prominent examples are the stars in a galaxy or in a globular cluster, and the dark matter particles in a galactic halo~\cite{BinneyTremaine-Book}. The state of such a gas at time $t$ is described by its one-particle distribution function (DF) $f(t,x,p)$, a nonnegative function on the one-particle phase space which determines the number of gas particles in a small phase space volume about the point $(x,p)$, and whose evolution is governed by the Vlasov equation.\footnote{The Vlasov equation is also often called the collisionless Boltzmann equation in the literature.} In the Newtonian gravitational case, the Vlasov equation is coupled to Poisson's equation for the gravitational potential generated by the gas, which yields the Vlasov-Poisson system, see Ref.~\cite{gR07} for a review of its mathematical properties. Its general relativistic counterpart is the Vlasov-Einstein system~\cite{hA11}.

It is instructive to compare the dynamics of a self-gravitating collisionless gas with the one of a pressureless fluid (dust), which is obtained in the limit in which the velocity dispersion of the gas vanishes. The gravitational collapse of a spherical dust cloud leads to the formation of a singularity in finite time at which the density diverges, both in the Newtonian theory and in general relativity~\cite{jOhS39}; in the latter case, part of this singularity may even be visible to distant observers~\cite{dC84,nOoS11}. In contrast to this, the velocity dispersion of a kinetic gas prevents the particles from focusing. In fact, the Vlasov-Poisson system admits global classical solutions for general initial data~\cite{kP92,pLbP91,jS91}, while for the spherically symmetric Vlasov-Einstein system global solutions exist for small initial data~\cite{gRaR92} and sufficiently concentrated data collapse to a black hole~\cite{hAmKgR11}. Therefore, for a Newtonian collisionless gas the relevant question is not whether or not a singularity forms, but rather how the gas behaves for large times and, in particular, whether or not it relaxes to a stationary state.

A fundamental mechanism leading to such a relaxation is phase space mixing~\cite{dL62,dL67}. To explain it, consider first a gas whose self-gravity is neglected and which is trapped in the potential well of an external potential. If the one-particle Hamiltonian is integrable, each gas particle is confined to an invariant torus, and in terms of action-angle variables $(Q,J)$ the solution of the Vlasov equation is $f(t,Q,J) = F(Q - \Omega(J)t,J)$, where $F$ describes the initial DF and $\Omega(J)$ denotes the frequencies associated with the torus labeled by $J$. Since these frequencies vary from one torus to the other, the DF is stretched along the angle directions and develops structures on finer and finer scales as time increases. As a consequence, $f$ does not converge pointwise as $t\to\infty$. However, the macroscopic observables which are obtained by integrating $f$ against a test function do converge, and their limit is the one associated with the average of the initial DF over the angle variables, which depends only on integrals of motion. Note that this holds for any initial DF, not only for small perturbations of a steady state. A mathematically precise formulation of this statement and its proof for rotationally symmetric potentials in $d = 1,2,3$ space dimensions satisfying a suitable non-degeneracy condition were given by Rioseco and Sarbach~\cite{pRoS20} (see also Refs.~\cite{sTmHdL86,cM19}). The results in Ref.~\cite{pRoS20} include the case of a spherically symmetric gas configuration in which all the gas particles have the same value of the total angular momentum, which behaves like an effective one-dimensional system, and they explain the relaxation to a virialized state which had been observed in the numerical simulations of Ref.~\cite{pDeJmAeMdN17}. The final state is thus determined by the initial DF through its average over the angle variables; for a gas in an axisymmetric potential with two potential wells along the symmetry axis, in which the motion is integrable, this prediction has recently been compared with $N$-particle simulations in Ref.~\cite{aGPmGdNoS26}. For the analogous phenomenon in the exterior of a black hole we refer the reader to Refs.~\cite{pRoS18,pRoS24}, and for estimates on the rate at which the observables converge to Refs.~\cite{mMpRhV22,sCjL22,mHgRmScS24}.

The generalization to the self-gravitating case, which was left open in Ref.~\cite{pRoS20}, is considerably more involved, since the potential now depends on the DF, such that a change of the latter induces a change of the potential which, in turn, acts back on the gas. In plasma physics, where the particles interact through their mean electric field, the corresponding collective effect has been known since the work of Landau~\cite{lL46}: for a gas of charged particles close to a spatially homogeneous equilibrium the electric field decays in time while the DF undergoes mixing. This phenomenon, known as Landau damping, has been established rigorously at the nonlinear level by Mouhot and Villani~\cite{cMcV11}. The gravitational case differs from the electromagnetic one in two important aspects. First, the interaction is attractive, such that the collective response of the gas tends to reinforce a perturbation instead of screening it. Second, a self-gravitating gas configuration of finite mass does not possess spatially homogeneous equilibria. Instead, its steady states are spatially inhomogeneous, the gas particles are trapped in the potential well, and the frequencies of their radial motion lie in an interval which, in the situations considered in these notes, is bounded and bounded away from zero. For this reason, a collective oscillation whose frequency lies below this interval is not in resonance with any gas particle and hence cannot be damped by mixing.

Two situations have been analyzed in the gravitational case. The first one concerns small perturbations of a self-consistent steady state. Had\v{z}i\'c, Rein and Straub~\cite{mHgRcS22} (see also Refs.~\cite{Kunze-Book,mK22}) showed that the essential spectrum of the operator governing the linearized dynamics about a large class of spherically symmetric steady states is determined by the radial periods of the gas particles, and, using a version of the Birman-Schwinger principle which was first applied to this problem by Mathur~\cite{sM90}, they derived a criterion for the existence of an eigenvalue below this spectrum, which corresponds to an oscillating mode that is not damped. Related criteria and estimates for the number of such modes have been given in Ref.~\cite{mMpRhV26}. For the reduced system in which all the gas particles have the same magnitude of the angular momentum and move around a central point mass, Had\v{z}i\'c, Rein, Schrecker and Straub~\cite{mHgRmScS25} proved the following dichotomy for steady states of sufficiently small mass whose DF is proportional to $(E_0 - E)^k$ for energies $E$ below the cutoff energy $E_0$: the solutions of the linearized system are damped if $k > 1$, whereas an oscillating mode exists if $1/2 < k\leq 1$. Quantitative decay estimates for the linearized system have been obtained in Refs.~\cite{mHgRmScS24,mHmS25}, and particle-in-cell simulations show that, depending on the steady state, its perturbations either oscillate without damping or decay~\cite{tRgR18,cS24}. The second situation concerns initial data which are not a perturbation of a steady state of the gas. For a self-gravitating gas around a point mass, Chaturvedi and Luk~\cite{sCjL26} proved that small, spherically symmetric initial data which are supported on bounded orbits, with angular momenta in an interval bounded away from zero, undergo phase space mixing up to a time of the order of $\epsilon^{-1}(\log\epsilon^{-1})^{-2}$, where $\epsilon$ measures the size of the data, while the global existence and the asymptotic behavior of these solutions remain open. Finally, when the initial data are far from any steady state and the self-gravity is strong, the gas undergoes a violent relaxation~\cite{dL67} toward a quasi-stationary state, and the prediction of this state from the initial data is understood only in special cases, see Ref.~\cite{yLrPfRtTfB14} for a review.

In these notes, we study numerically the relaxation through phase space mixing of a spherically symmetric collisionless gas which is trapped in an external central potential, whose self-gravity is taken into account, and whose initial DF is not a steady state. We consider the reduced model in which all the gas particles have the same magnitude $L_0 > 0$ of the angular momentum, which has a two-dimensional phase space and is the setting considered in Refs.~\cite{pRoS20,mHgRmScS25}. This model goes back to the numerical experiments of H\'enon with spherical shells (see Ref.~\cite{mH73} and references therein) and has been used in the $N$-body simulations of Ref.~\cite{pKbM02}. For the external potential we choose the isochrone potential, for which the action-angle variables are known explicitly, and the potential of a point mass, which is the one considered in Refs.~\cite{mHgRmScS25,sCjL26}. We address three questions: how fast the macroscopic observables relax and what determines the rate, what can prevent the mixing from being completed, and which state the gas reaches at late times. Our analysis is based on three tools: (i) a particle-in-cell (PIC) code which solves the nonlinear Vlasov-Poisson system, (ii) a numerical solution of the equations linearized about a self-consistent steady state, formulated in the action-angle variables of the latter, which describes initial data close to that steady state, and (iii) the numerical computation of the largest eigenvalue $\lambda(\omega)$ of the operator of Mathur, which measures the gain of the self-consistent loop in which a perturbation of the potential with frequency $\omega$ induces a perturbation of the density which, in turn, generates a potential. An oscillating mode with frequency below the frequency band of the gas particles exists if and only if the value $\lambda_{edge}$ of this gain at the lower edge of the band is larger than one. This is the criterion of Refs.~\cite{sM90,mHgRcS22,mK22}; the quantity $\lambda_{edge}$ corresponds to the number which is denoted by $M$ in Ref.~\cite{mHgRcS22} and by $\mu_*$ in Ref.~\cite{mK22}.

The main results of these notes can be summarized as follows. First, in the absence of self-gravity the observables relax algebraically, with an exponent which is determined by the behavior of the initial DF at the boundary of its support: a DF which vanishes like $(J_{max} - J)^k$ at the outer edge of its support gives rise to a tail proportional to $t^{-(k+1)}$, and the PIC code reproduces the exact solution of this problem. Second, for initial data close to a self-consistent steady state, the self-gravity of the deviation from that state modifies the relaxation in two ways. If $\lambda_{edge} < 1$, the deviation is damped, first exponentially, as in Landau damping, and later algebraically; at late times, or for larger masses, the self-gravity reduces the exponent of the algebraic decay from $k + 1$ to $k$, which is the exponent of the decay estimates of Ref.~\cite{mHgRmScS24} for perturbations that vanish at the edge like the derivative of the DF of the steady state. If $\lambda_{edge} > 1$, the mixing is not completed, since the part of the deviation which is carried by the oscillating mode persists. For the families of steady states we consider, this happens above a threshold mass, and the frequency predicted by the equation $\lambda(\omega) = 1$ agrees with the one measured in the simulations. Third, the oscillating modes which lie close to the edge of the band can only be sustained with a small amplitude: above a critical amplitude, which we estimate by means of a pendulum model, the gas particles near the edge of the steady state are trapped by the mode. That the amplitude of such modes is limited was observed by Weinberg~\cite{mW91} in simulations of the one-dimensional slab; the pendulum model provides an explicit estimate for the amplitude at which this limitation sets in. Fourth, for initial data far from any steady state, in which the gas is concentrated about one orbital phase, the outcome is determined by the ratio $\mu$ between the bounce frequency of the gas particles in the potential well of the gas and the spread of their orbital frequencies: for small values of $\mu$ the gas mixes, whereas above a critical value of $\mu$, which lies between $0.3$ and $0.5$ in our simulations, it collapses into a self-bound core which does not mix.\EJ{Preliminary result from the runs B0--B5 of the demo document. To be consolidated in section~\ref{Sec:FarFromEq}: more masses, the dependence on the widths of the initial DF, convergence, and the final state.}

Particle methods have a long tradition in the numerical study of spherically symmetric collisionless systems, both in the Newtonian case~\cite{tRgR18,cS24} and in the relativistic one, where they have been used to analyze the collapse of star clusters and the critical behavior at the threshold of black hole formation~\cite{sStT85,gRaRjS98,Olabarrieta:2000bn,Olabarrieta:2001wy,hAgR06}; see Refs.~\cite{BirdsallLangdon-Book,HockneyEastwood-Book,lapenta2015kinetic} for introductions to the PIC method, Ref.~\cite{Shalaby:2017yok} for a higher-order implementation in the context of plasma physics, and Ref.~\cite{gRtR03} for a convergence proof of a PIC scheme for the spherically symmetric Vlasov-Einstein system. Measuring the relaxation of the observables with such a method is delicate, because the signal has to be separated from the noise which is due to the finite number of computational particles. In these notes we address this problem by placing the particles on a regular lattice in action-angle variables, by encoding a deviation from a steady state in the weights of the particles instead of their positions (as in Refs.~\cite{fLfCjB93,jBaOyY11}), and by subtracting a reference run without this deviation. A loading of the particles which reduces the initial fluctuations (a quiet start) was already used for systems of spherical shells in Ref.~\cite{mH73}. We also show that the observables need to be computed with the action-angle variables belonging to the total potential, since those of the external potential alone leave a time-independent residue which could be mistaken for an incomplete mixing. For initial data far from equilibrium there is no fixed potential with respect to which the gas relaxes, and we use the action-angle variables of the instantaneous potential.

The remainder of these notes is organized as follows. In section~\ref{Sec:VPSphSym} we formulate the Vlasov-Poisson system with an external potential in spherical symmetry and the reduced model with fixed angular momentum, and we introduce the action-angle variables and the observables used throughout these notes. In section~\ref{Sec:Linear} we review the mixing in the absence of self-gravity, derive the linearized equations about a steady state and state the criterion for the existence of oscillating modes. Our numerical methods are described in section~\ref{Sec:Numerics}, and their validation is presented in section~\ref{Sec:Validation}. Section~\ref{Sec:FixedL} is devoted to initial data close to a self-consistent steady state, and section~\ref{Sec:FarFromEq} to initial data far from equilibrium. Conclusions and an outlook, including the extension to a distribution of angular momenta, are given in section~\ref{Sec:Conclusions}. Technical details regarding the numerical computation of the action-angle variables are included in Appendix~\ref{App:ActionAngle}.
